\documentclass[11pt,a4paper]{article}
\usepackage[utf8]{inputenc}
\usepackage{graphicx}
\usepackage[left=2.5cm,top=3cm,right=2.5cm,bottom=3cm,bindingoffset=0.5cm]{geometry}
\usepackage{AEDLogica, AEDEspecificacion, AEDTADs}
\usepackage{caratula}
\newlength{\miSangria}
\setlength{\miSangria}{2em}

% Asi pueden escribir nuevos comandos. 
% Este por ejemplo asegura q los nombres 
% que figuren con una tipografia diferenciada  
\newcommand{\Tipo}[1]{\mathsf{#1}}
% la sintaxis es \newcommand{\nombreDeLaMacro}[cantidadDeParametros]{Lo que va ser remplazado por el macro} 
\newcommand{\norm}[1]{\vert#1\vert}


\begin{document}

\section{Titulo del Proc ACA}

\noindent topeDePromocion \textbf{ES}: $\mathbb{R}$
\vspace{0.5cm}

\noindent enum Categoría \textbf{ES} \{Bronce, Plata,Oro, Platino  \}

\vspace{0.5cm}

\noindent Transacciones \textbf{ES} Struct$\struct{
		UsuarioOrigen: \Tipo{\Z}, \\
		UsuarioDestino: \Tipo{\Z},
		TipoDeTransaccion: \Tipo{\Z},
		Millas: \Tipo{\R}, \\
		NombrePromo: \Tipo{\str},
		FactorPromo: \Tipo{\R}
	}$

\begin{tad}{Airmiles}


	\obs{Usuarios}{\dict{id:\Tipo{\Z}}{CategoriaActual : Categoria}}

	\noindent \obs{HistorialDeTransacciones}{\seq{Transacciones}}
	\\
	\texttt{obs} Promociones : dict \{ \\
	\hspace*{1.5em} NombrePromocion: $\str$, \\
	\hspace*{1.5em} $\struct{FactorDePromocion:\Tipo{\R},TopeDePromocion:\Tipo{\R}, MillasRestantes:\Tipo{\R}}$ \\
	\}
	\vspace{0.5cm}


	\begin{proc}{ResumenTransacciones}
		{\In AM: Airmillas, \In id : \Tipo{\Z}}
		{\tupla{\Tipo{\R}, \Tipo{\R}, \Tipo{\R}}}
		\requiereLargo{ id \in claves(AM.Usuarios)}

		\aseguraLargo{ res_{0} = MillasObtenidasPorVuelo(AM.Transacciones, id) \yLuego
		\\
		res_{1} = MillasDescontadasPorCanjeo(AM.Transacciones, id) \yLuego
		\\
		res_{2} = MillasTransferidas(AM.Transacciones, id, id)
		}
	\end{proc}


	\auxLargo{MillasObtenidasPorVuelo}
	{Historial: \Tipo{\seq{Transacciones}}, idOrigen: \Tipo{\Z}}
	{\Tipo{\R}}
	{
		\sum\limits_{i=0}^{\norm{Historial} - 1}
		\IfThenElse{
			EsVuelo(Historial[i],
			idOrigen)
		}{
			Historial[i].Millas
		}{
			0
		}
	}

	\vspace{1cm}

	\auxLargo{MillasTransferidas}
	{Historial: \Tipo{\seq{Transacciones}}, idOrigen: \Tipo{\Z}, idDestino: \Tipo{\Z}}
	{\Tipo{\R}}
	{
		\sum\limits_{i=0}^{\norm{Historial} - 1}
		\IfThenElse{
			EsTransferencia(Historial[i],
			\\
			idOrigen, idDestino)
		}{
			Historial[i].Millas
		}{
			0
		}
	}

	\vspace{1cm}
	\auxLargo{MillasDescontadasPorCanjeo}
	{Historial: \Tipo{\seq{Transacciones}}, idOrigen: \Tipo{\Z}}
	{\Tipo{\R}}
	{
		\sum\limits_{i=0}^{\norm{Historial} - 1}
		\IfThenElse{
			EsCanjeo(Historial[i],
			idOrigen)
		}{
			Historial[i].Millas
		}{
			0
		}
	}
    \vspace{1cm}
    \predLargo{EsVuelo}
        {Transaccion: \Tipo{Struct}, idOrigen: \Tipo{\Z}}
        {
            Transaccion.UsuarioOrigen = idOrigen \land
            \\
            Transaccion.TipoDeTransaccion = 0 \land
            \\
            Transaccion.UsuarioDestino = -1 \land
            \\
            \comentario{Deberia verificar los otros campos? Se complica...}
        }

    \vspace{1cm}

    \predLargo{EsTransferencia}
        {Transaccion: \Tipo{struct}, idOrigen: \Tipo{\Z}, idDestino: \Tipo{\Z}}
        {
            Transaccion.UsuarioOrigen = idOrigen \land
            \\
            Transaccion.TipoDeTransaccion = 1 \land
            \\
            Transaccion.UsuarioDestino = idDestino \land
            \\
            Transaccion.millas > 0
            \\
            \comentario{Deberia verificar los otros campos? Se complica...}
        }

    \vspace{1cm}

    

    \predLargo{EsCanjeo}
        {Transaccion: \Tipo{struct}, idOrigen: \Tipo{\Z}}
        {
            Transaccion.UsuarioOrigen = idOrigen \land
            \\
            Transaccion.TipoDeTransaccion = 2 \land
            \\
            Transaccion.UsuarioDestino = -1 \land
            \\
            Transaccion.distancia > 0
            \\
            \comentario{Deberia verificar los otros campos? Se complica...}
        }



\end{tad}
\end{document}